% !TEX program = pdflatex
\documentclass[11pt,a4paper]{article}

% =========================================================================
% FICHA TÉCNICA DE LA PROPUESTA
% Proyecto: Horus Sentinel
% Gestión de Proyectos TICS (EST315) -- Semestre 2026-II -- Team Sobrevivientes
% =========================================================================

% --- PAQUETES DE CONFIGURACIÓN ---
\usepackage[utf8]{inputenc}
\usepackage[T1]{fontenc}
\usepackage[spanish,es-tabla]{babel}
% Compatibilidad con babel-spanish reciente (flechas de TikZ y signo de porcentaje)
\AtBeginDocument{\shorthandoff{<>}\spanishplainpercent}
\usepackage{geometry}
\geometry{
    top=3.6cm,
    bottom=2.3cm,
    left=2.2cm,
    right=2.2cm,
    % PARTE MODIFICABLE: headheight debe ser mayor que la altura del logo más alto (2.56cm = 73pt).
    headheight=80pt,
    headsep=15pt,
    footskip=20pt
}

\usepackage{graphicx}
\usepackage{fancyhdr}
\usepackage[table]{xcolor}
\usepackage{titlesec}
\usepackage{booktabs}
\usepackage{tabularx}
\usepackage{array}
\usepackage{amsmath}
\usepackage{enumitem}
\usepackage{microtype}
\usepackage{hyperref}

% Logos: funciona a 2 o 3 niveles de profundidad dentro de docs/latex/
\graphicspath{{../../assets/img/}{../../../assets/img/}}

% --- PALETA INSTITUCIONAL FIEI / UNA PUNO ---
\definecolor{AzulUNA}{RGB}{10, 45, 90}
\definecolor{DoradoFIEI}{RGB}{180, 130, 20}
\definecolor{GrisOscuro}{RGB}{50, 50, 50}
\definecolor{AzulClaro}{RGB}{228, 236, 246}

% --- CONFIGURACIÓN DE LISTAS ---
\setlist{topsep=2pt, itemsep=2pt, parsep=1pt}

% --- COLUMNAS DE TABLA ---
\newcolumntype{L}{>{\raggedright\arraybackslash}X}
\newcolumntype{P}[1]{>{\raggedright\arraybackslash}p{#1}}
\newcolumntype{C}[1]{>{\centering\arraybackslash}p{#1}}
\renewcommand{\arraystretch}{1.25}

% --- FORMATO DE ENLACES ---
\hypersetup{
    colorlinks=true,
    linkcolor=AzulUNA,
    urlcolor=AzulUNA,
    citecolor=AzulUNA,
    pdftitle={Ficha Técnica de la Propuesta -- Horus Sentinel},
    pdfauthor={Team Sobrevivientes}
}

% --- FORMATO DE TÍTULOS ---
\titleformat{\section}{\color{AzulUNA}\normalfont\large\bfseries}{\thesection.}{0.7em}{}[\color{AzulUNA}\titlerule]
\titleformat{\subsection}{\color{AzulUNA}\normalfont\normalsize\bfseries}{\thesubsection.}{0.5em}{}
\titlespacing*{\section}{0pt}{9pt}{4pt}
\titlespacing*{\subsection}{0pt}{6pt}{2pt}

% --- ENCABEZADO Y PIE DE PÁGINA (fancyhdr) ---
\pagestyle{fancy}
\fancyhf{}

% ===== TAMAÑO DE LOS LOGOS (PARTE MODIFICABLE) =====
% logo_fiei.png  : height=1.45cm (lado izquierdo)
% logo_equipo.png: height=2.56cm (lado derecho; 25 % más grande porque la imagen tiene más borde blanco)
% Para cambiar el tamaño, modifique solo el valor height=... de cada \includegraphics.
% Si el logo más alto supera 2.6cm, aumente también headheight y top en la configuración de página.
\fancyhead[L]{%
    \includegraphics[height=1.45cm,keepaspectratio]{logo_fiei.png}%
}
\fancyhead[C]{%
    \scriptsize \textbf{\textcolor{AzulUNA}{UNIVERSIDAD NACIONAL DEL ALTIPLANO DE PUNO}}\\[2pt]
    \tiny \textcolor{GrisOscuro}{FACULTAD DE INGENIERÍA ESTADÍSTICA E INFORMÁTICA}\\[2pt]
    \scriptsize \textbf{\textcolor{AzulUNA}{Gestión de Proyectos TICS (EST315) \textbar\ Team Sobrevivientes}}\\[2pt]
    \footnotesize \textbf{\textcolor{DoradoFIEI}{FICHA TÉCNICA DE LA PROPUESTA}}%
}
\fancyhead[R]{%
    \includegraphics[height=2.56cm,keepaspectratio]{logo_equipo.png}%
}

\renewcommand{\headrulewidth}{0pt}
\renewcommand{\headrule}{}

\fancyfoot[L]{\scriptsize \textcolor{GrisOscuro}{Horus Sentinel}}
\fancyfoot[C]{\small \thepage}
\fancyfoot[R]{\scriptsize \textcolor{GrisOscuro}{Puno, Perú -- Semestre 2026-II}}
\renewcommand{\footrulewidth}{0.4pt}

\begin{document}

\begin{center}
    {\Large \textbf{\textcolor{AzulUNA}{FICHA TÉCNICA DE LA PROPUESTA}}}\\[3pt]
    {\normalsize \textbf{Horus Sentinel -- Sistema de mapeo predictivo de zonas de riesgo delictivo}}\\[4pt]
    \footnotesize \textbf{Fecha:} miércoles 02/09/2026 (semana 2) \quad \textbar\quad \textbf{Versión:} 1.0 \quad \textbar\quad \textbf{Fase:} Inicio
\end{center}

\vspace{0.1cm}

\section{Datos Generales}
\noindent\begin{tabularx}{\textwidth}{@{}P{4.2cm} L@{}}
\toprule
\textbf{Nombre del proyecto:} & Horus Sentinel \\
\textbf{Tipo de solución:} & Plataforma geoespacial con mapa de calor de riesgo y sugerencia de ruta segura \\
\textbf{Enfoque principal:} & Estadística espacial para estimar dónde y cuándo es más probable un incidente \\
\textbf{Plataforma:} & Aplicación con mapa interactivo \\
\textbf{Ámbito:} & Un distrito piloto de la ciudad de Puno \\
\textbf{Complejidad estimada:} & Media a alta: modelado estadístico espaciotemporal y geolocalización \\
\textbf{Enfoque metodológico:} & Ágil (Scrum), por fases: recolectar datos, construir el modelo, armar el mapa y el ruteo, probar en la zona piloto y ajustar \\
\textbf{Equipo:} & Team Sobrevivientes \\
\bottomrule
\end{tabularx}

\section{Problema u Oportunidad}
No existe información clara y actualizada sobre qué zonas de la ciudad son más peligrosas y en qué horarios. Por eso, el patrullaje y los desplazamientos de las personas se deciden sin datos que los respalden.

\section{Objetivo General}
Desarrollar un sistema que analice reportes de incidentes y muestre un mapa de calor de riesgo por zona y horario, útil para orientar el patrullaje y para que ciudadanos y turistas elijan rutas más seguras.

\section{Objetivos Específicos}
\begin{enumerate}[leftmargin=*]
    \item Analizar los datos disponibles de incidentes para estimar el riesgo por zona y horario.
    \item Detectar zonas con concentración inusual o creciente de incidentes.
    \item Mostrar los resultados en un mapa interactivo con sugerencia de rutas seguras.
\end{enumerate}

\section{Beneficiarios}
{\small
\noindent\begin{tabularx}{\textwidth}{@{}P{4.2cm} L@{}}
\toprule
\rowcolor{AzulClaro}
\textbf{Beneficiario} & \textbf{Beneficio} \\
\midrule
Ciudadanos y estudiantes & Eligen rutas y horarios de desplazamiento con menor riesgo. \\
Turistas & Conocen las zonas de mayor riesgo antes de recorrer la ciudad. \\
Serenazgo y Policía & Priorizan el patrullaje con información por zona y horario. \\
Municipalidad & Orienta inversiones en alumbrado y vigilancia. \\
\bottomrule
\end{tabularx}
}

\section{Alcance}
{\small
\noindent\begin{tabularx}{\textwidth}{@{}L L@{}}
\toprule
\rowcolor{AzulClaro}
\textbf{Incluye} & \textbf{No incluye} \\
\midrule
Mapa de riesgo por zona y horario & Cobertura nacional (se inicia con un distrito piloto) \\
Sugerencia de la ruta más segura & Identificación de personas o sospechosos \\
Detección de zonas con patrones anómalos & Cámaras o vigilancia en tiempo real \\
\bottomrule
\end{tabularx}
}

\section{Justificación}
La inseguridad ciudadana es un problema visible y los recursos de patrullaje son limitados. Usar datos para priorizar dónde y cuándo actuar es más eficiente que decidir sin información.

\section{Técnicas Principales}
\begin{itemize}[leftmargin=*]
    \item Estimación de densidad por kernel (KDE) para el mapa de calor base.
    \item Proceso de Poisson no homogéneo espaciotemporal para modelar cómo un incidente eleva la probabilidad de otro cercano.
    \item Algoritmo $A^*$ sobre el grafo de calles con costo $= \alpha \cdot \text{distancia} + \beta \cdot \text{riesgo}$.
\end{itemize}

\section{Principal Valor}
Convierte datos dispersos de criminalidad en información preventiva y útil para decidir, como el ojo que todo lo ve de Horus.

\section{Riesgo Clave}
Sesgo hacia zonas históricamente más vigiladas, que podría concentrar el patrullaje donde más se denuncia y no donde más ocurre. Respuesta: combinar fuentes oficiales con reportes ciudadanos y auditar el modelo.

\section{Firmas}
\noindent En señal de conformidad con la presente propuesta, los integrantes del equipo firman a continuación:

\vspace{1.5cm}
\begin{center}
\begin{tabular}{c c c}
    \parbox{6.5cm}{\centering \rule{6.0cm}{0.6pt}\\[5pt] \textbf{Clyde Neil Paricahua Pari}\\[2pt] {\footnotesize Project Manager}}
    & \hspace{0.8cm} &
    \parbox{6.5cm}{\centering \rule{6.0cm}{0.6pt}\\[5pt] \textbf{Ruth Karina Apaza Solis}\\[2pt] {\footnotesize Analista de Planificación}}
    \\[2.0cm]
    \parbox{6.5cm}{\centering \rule{6.0cm}{0.6pt}\\[5pt] \textbf{Eliseo Tarqui Ajahuana}\\[2pt] {\footnotesize Especialista Técnico / Arquitectura}}
    & \hspace{0.8cm} &
    \parbox{6.5cm}{\centering \rule{6.0cm}{0.6pt}\\[5pt] \textbf{Viviana Mary Quisbert Quispe}\\[2pt] {\footnotesize Gestora de Comunicación, Riesgos y Documentación}}
\end{tabular}
\end{center}
\end{document}
